% =====================================================================
%  TYPE: Postoperative instructions -- lumbar (percutaneous decompression, discectomy, fusion)
%  Sections: patient block, instructions, follow-up, emergency
%  box, when to call, medication page. Based on Kokkonouzi / Zampetakis 2026-04.
%
%  Copy, rename to "YYYYMMDD - surname - topic.tex", fill in and
%  compile with XeLaTeX. Toggle alternatives by moving the "%".
%  Text in [brackets] is a placeholder: replace or delete it.
% =====================================================================
\documentclass[10pt,greek]{article}
\usepackage[useregional]{datetime2}

%%%%%%%%%%%%%%%%%%%%% Language
\providecommand\DocLang{gr}

%%%%%%%%%%%%%%%%%%%%% Clinic (letterhead logo, footer address)
\providecommand\Clinic{Peristeri}
% \providecommand\Clinic{Mandra}
% \providecommand\Clinic{Bioclinic}
% \providecommand\Clinic{Psychikou}
% \providecommand\Clinic{Metropolitan}
% \providecommand\Clinic{Dunant}
% \providecommand\Clinic{Euroclinic}

%%%%%%%%%%%%%%%%%%%%% Document code for the QR: paste the code issued
%%%%%%%%%%%%%%%%%%%%% by doc.erns.eu (empty = no QR code)
\providecommand\DocCode{}

%%%%%%%%%%%%%%%%%%%%% Patient
\providecommand\PatientPrefix{Ο}
% \providecommand\PatientPrefix{Η κα}
\providecommand\PatientName{Παπαδόπουλος Γιώργος}
\providecommand\PatientAge{45}
% \providecommand\PatientAge{}

%%%%%%%%%%%%%%%%%%%%% Date
\providecommand\DocDate{\DTMdate{2026-01-15}}
% \providecommand\DocDate{\DTMdate{2026-05-04}}

%%%%%%%%%%%%%%%%%%%%% Title
\providecommand\DocTitle{Μετεγχειρητικές Οδηγίες}
% \providecommand\DocTitle{Ενημερωτικό Σημείωμα για τον Ασθενή}

%%%%%%%%%%%%%%%%%%%%% Closing (empty = no signature block)
\providecommand\DocClosing{Ο ιατρός}
% \providecommand\DocClosing{}

%%%%%%%%%%%%%%%%%%%%% Stamp and signature image
\providecommand\DocStamp{no}
% \providecommand\DocStamp{yes}

\input{evrtemplate}

\begin{document}

\begin{PatientInfo}
  \Info{Διάγνωση}{Τιμή 1}
  \Info{Χειρουργείο}{(10/01/2026) Τιμή 2}
%   \Info{Προηγηθέν}{Τιμή 4}
%   \Info{Συνοδά}{Τιμή 5}
\end{PatientInfo}

Ακολουθούν οδηγίες για την ομαλή σας ανάρρωση:

\section*{Γενικές Οδηγίες}
\begin{enumerate}[label=\arabic*.]
\setlength{\parskip}{6pt}
  \item \textbf{Κινητοποίηση}
       \begin{itemize}
% \item Παραμείνετε κλινήρης για Τιμή 6 ημέρες, όπως σας υποδείχθηκε από τον θεράποντα ιατρό.
% \item Επιτρέπεται σύντομη μετακίνηση μόνο για φυσικές ανάγκες και πάντα με βοήθεια.
\item Χρησιμοποιείτε \textbf{πάντα} τη ζώνη οσφύος κατά την όρθια στάση ή βάδιση.
\item Αποφύγετε για τουλάχιστον 6 εβδομάδες: κάμψεις, στροφές κορμού και άρση βάρους.
\item Μετά την επανεξέταση προτείνεται ήπιο περπάτημα με σταδιακή αύξηση.
    \end{itemize}
  \item \textbf{Πόνος -- Φάρμακα}
      \begin{itemize}
\item Λαμβάνετε τα παυσίπονα όπως έχουν συνταγογραφηθεί.
\item Επιδείνωση ή αλλαγή του χαρακτήρα του πόνου πρέπει να αναφέρεται άμεσα.
% \item \textbf{Μην λαμβάνετε αντιφλεγμονώδη} (π.χ.\ Voltaren, Ponstan, Naprosyn, Brufen) χωρίς ιατρική οδηγία.
    \end{itemize}
  \item \textbf{Φροντίδα Τραύματος}
      \begin{itemize}
\item Το τραύμα πρέπει να παραμένει καθαρό και στεγνό.
\item Αλλαγή τραύματος κάθε 2 ημέρες με αυτοκόλλητα γαζάκια για 5 ημέρες.
\item Αποφύγετε διαβροχή της περιοχής όσο υπάρχουν γάζες.
\item Μετά τις 5 ημέρες μπορεί να χρησιμοποιείται spray τύπου Pulvo47 ή άλλο αντίστοιχο, 2 φορές την ημέρα.
% \item Μην εφαρμόζετε κρέμες ή αλοιφές χωρίς οδηγία.
% Διαρροή ΕΝΥ: \item Αλλαγή τραύματος κάθε 2 ημέρες με αυτοκόλλητα γαζάκια για 10 ημέρες.
    \end{itemize}
  \item \textbf{Διατροφή}
    \begin{itemize}
\item Ελαφριά διατροφή, πλούσια σε πρωτεΐνες και φυτικές ίνες.
\item Επαρκής ενυδάτωση ημερησίως.
\item Σε περίπτωση δυσκοιλιότητας, μπορεί να χρησιμοποιηθεί ήπιο καθαρτικό (π.χ.\ nujol), εφόσον χρειάζεται.
    \end{itemize}
  \item \textbf{Οδήγηση}

Αποφύγετε την οδήγηση μέχρι νεότερη ιατρική οδηγία. Η οδήγηση μπορεί να επηρεαστεί από τον μετεγχειρητικό πόνο, τη δυσκαμψία στη μέση, την περιορισμένη ικανότητα απότομων κινήσεων και τα παυσίπονα/αντιφλεγμονώδη που ενδέχεται να προκαλούν υπνηλία ή καθυστέρηση αντίδρασης.
% Επιτρέπεται να οδηγήσετε μόνο όταν:
%   \begin{itemize}[nosep]
% \item κινείστε άνετα χωρίς πόνο κατά το φρενάρισμα ή την απότομη στροφή,
% \item δεν λαμβάνετε φάρμακα που επηρεάζουν την εγρήγορση (π.χ.\ οπιοειδή),
% \item μπορείτε να καθίσετε χωρίς δυσφορία για 20--30 λεπτά.
%   \end{itemize}
% \item \textbf{Σεξουαλική δραστηριότητα:} Εφόσον δεν υπάρχει πόνος, μετά τις 4--6 εβδομάδες.
\end{enumerate}

\section*{Εξειδικευμένες Οδηγίες}
\begin{enumerate}[label=\arabic*.]
  \item \textbf{Ζώνη οσφύος}
      \begin{itemize}
\item Χρήση μόνο κατά την κινητοποίηση.
\item Δεν απαιτείται κατά την κατάκλιση ή τον ύπνο.
    \end{itemize}
  \item \textbf{Φυσικοθεραπεία}
      \begin{itemize}
\item Δεν απαιτείται άμεσα μετά το χειρουργείο.
\item Τον πρώτο μήνα αρκεί ήπια δραστηριότητα και περπάτημα.
% \item Ασκήσεις ενδυνάμωσης του πρόσθιου κνημιαίου μυός σε ξαπλωτή θέση για τον πρώτο μήνα.
% \item Θα πρέπει να ξεκινήσετε άμεσα συνεδρίες εργοθεραπείας και κινησιοθεραπείας με μαλάξεις και ηλεκτρική διέγερση για την επαναφορά της κινητικότητας στο κάτω άκρο.
\item Μετά την πρώτη επανεξέταση θα δοθούν περαιτέρω οδηγίες, εφόσον χρειάζεται.
    \end{itemize}
% \item \textbf{Επιστροφή στην εργασία:} Μετά από Τιμή 7 εβδομάδες, ανάλογα με το είδος της εργασίας.
  \item \textbf{Αντιθρομβωτική αγωγή}
      \begin{itemize}
\item Κάθε ημέρα Τιμή 8 υποδορίως, για Τιμή 9 ημέρες, την ίδια ώρα.
\item Αποφύγετε μάλαξη στο σημείο ένεσης.
% \item Η τακτική ήπια βάδιση συμβάλλει στην πρόληψη θρόμβωσης.
    \end{itemize}
    \item \textbf{Μπάνιο -- Υγιεινή}
        \begin{itemize}
\item Μπάνιο επιτρέπεται με χρήση αδιάβροχων αυτοκόλλητων γαζών. Αποφύγετε τη διαβροχή του τραύματος. Εάν αυτό συμβεί, στεγνώστε καλά με αποστειρωμένες γάζες και καλύψτε με αυτοκόλλητο.
\item Ολοκληρωτικό ντους επιτρέπεται μετά από 10 ημέρες και μετά από σύσταση του ιατρού.
% \item \textbf{Απαγορεύεται} το μπάνιο σε μπανιέρα, πισίνα ή θάλασσα.
        \end{itemize}
\end{enumerate}

\section*{Επανέλεγχοι}
\begin{description}
\item[Σε 10 ημέρες:] κλινική εξέταση και αφαίρεση ραμμάτων.
\item[Σε 1 μήνα:] επανεκτίμηση και οδηγίες για δραστηριότητες / φυσικοθεραπεία.
% \item[Σε 3 μήνες:] αξονική τομογραφία ΟΜΣΣ.
\end{description}

% Emergency box: not in the old lumbar letters. The items are taken from
% the list below; review before enabling.
% \begin{UrgentBox}
% \item \textbf{Ακράτεια ούρων/κοπράνων} ή \textbf{αδυναμία ούρησης}.
% \item \textbf{Νέα ή ταχέως επιδεινούμενη αδυναμία} στα κάτω άκρα.
% \end{UrgentBox}

\textbf{Επικοινωνήστε άμεσα με τον ιατρό εάν εμφανίσετε:}
\begin{itemize}
\item Έντονο ή επιδεινούμενο πόνο.
\item Πυρετό (Θερμοκρασία $>$ 38\textdegree{}C).
\item Ερυθρότητα, οίδημα ή εκροή από το τραύμα.
\item Αδυναμία ή αιμωδίες στα κάτω άκρα.
\item Ακράτεια ούρων/κοπράνων ή αδυναμία ούρησης.
\end{itemize}

\newpage
\begin{center}
\LARGE\textbf{Οδηγίες Φαρμακευτικής Αγωγής}
\end{center}
\vspace{0.5cm}

\begin{MedTable}
\DrugEntry{inj Clexane}{4000IU}{1 x 1}{10}{αντιπηκτικές ενέσεις υποδορίως, την ίδια ώρα κάθε ημέρα}
\DrugEntry{tb Zinadol}{500mg}{1 x 2}{5}{αντιβίωση}
\DrugEntry{tb Movatec}{7.5mg}{1 x 2}{10}{μετά το φαγητό}
\DrugEntry{tb Apotel}{1000mg}{2 ή 3 την ημέρα}{}{παυσίπονο, μέγιστη δόση 3 δισκία το 24ωρο}
\DrugEntry{caps Omeprazole}{40mg}{1 x 2}{10}{προ φαγητού}
% ---- Drug library: uncomment the lines you need ---------------------
% -- Analgesics
% \DrugEntry{tb Paracetamol}{1000mg}{1 x 3}{10}{}
% \DrugEntry{tb Apotel}{1000mg}{2 ή 3 την ημέρα}{}{παυσίπονο, μέγιστη δόση 3 δισκία το 24ωρο}
% \DrugEntry{tb Depon}{500mg}{1 ή 2, έως 3 φορές την ημέρα}{}{εφόσον υπάρχει πόνος -- μέγιστο 3g το 24ωρο}
% \DrugEntry{tb Tramal}{50mg}{1 x 2 ή 3}{6}{}
% \DrugEntry{tb Zaldiar}{}{2 ή 3 την ημέρα}{}{εφόσον υπάρχει έντονος πόνος}
% -- NSAIDs
% \DrugEntry{tb Movatec}{7.5mg}{1 x 2}{10}{μετά το φαγητό}
% \DrugEntry{tb Movatec}{15mg}{1 x 1}{1}{μετά το φαγητό}
% \DrugEntry{tb Naprosyn}{500mg}{1 x 2}{6}{μετά το φαγητό}
% -- Steroids (tapers)
% \DrugEntry{tb Medrol}{4mg}{1 x 3\newline 1 x 2\newline 1 x 1\newline STOP}{}{για 5 ημέρες\newline για 5 ημέρες\newline για 5 ημέρες}
% \DrugEntry{tb Medrol}{4mg}{2 x 3\newline 1 x 3\newline 1 x 2\newline 1 x 1\newline STOP}{}{για 3 ημέρες\newline για 3 ημέρες\newline για 3 ημέρες\newline για 3 ημέρες}
% \DrugEntry{syr Dexaton}{2mg/5ml}{10cc x 3}{10}{μετά το φαγητό}
% \DrugEntry{tb Dexamethasone}{4mg}{1 x 3}{}{μετά το φαγητό}
% -- Antibiotics
% \DrugEntry{tb Zinadol}{500mg}{1 x 2}{5}{αντιβίωση}
% \DrugEntry{tb Klaricid}{500mg}{1 x 2}{7}{αντιβίωση}
% \DrugEntry{caps Amoxicillin}{500mg}{1 x 3}{7}{μετά το φαγητό}
% -- Gastroprotection
% \DrugEntry{caps Omeprazole}{40mg}{1 x 2}{10}{προ φαγητού}
% \DrugEntry{caps Omeprazole}{20mg}{1 x 2}{10}{15 λεπτά πριν το φαγητό}
% \DrugEntry{caps Mostrelan}{40mg}{1 x 1}{10}{προ φαγητού}
% -- Anticoagulants
% \DrugEntry{inj Clexane}{4000IU}{1 x 1}{10}{αντιπηκτικές ενέσεις υποδορίως, την ίδια ώρα κάθε ημέρα}
% \DrugEntry{inj Innohep}{0,45ml}{1 x 1}{10}{αντιπηκτικές ενέσεις υποδορίως, την ίδια ώρα κάθε ημέρα}
% \DrugEntry{inj Fraxiparine}{2850IU}{1 x 1}{10}{αντιπηκτικές ενέσεις υποδορίως, την ίδια ώρα κάθε ημέρα}
% -- Neuropathic pain, antiepileptics
% \DrugEntry{tb Lyrica}{50mg}{1 x 1 (βράδυ)}{4}{μετά την 4η ημέρα, 1 x 2}
% \DrugEntry{caps Lyrica}{150mg}{1 x 2}{}{πρωί -- βράδυ, με τροφή. Μην διακόπτεται απότομα}
% \DrugEntry{tb Neurontin}{400mg}{2 x 2}{}{}
% \DrugEntry{tb Tegretol}{200mg}{1 x 3}{}{με τροφή, σταθερές ώρες. \textbf{Μην διακόπτεται απότομα}}
% \DrugEntry{caps Cymbalta}{60mg}{1 x 2}{}{}
% \DrugEntry{tb Keppra}{1000mg}{1 x 3}{}{}
% \DrugEntry{tb Depakine}{400mg}{1 x 3}{}{θα γίνει προοδευτική μείωση σε συνεννόηση με τον θεράποντα ιατρό}
% -- Other
% \DrugEntry{sir Duphalac}{5ml}{1 ή 2 την ημέρα}{}{καθαρτικό, ανάλογα με τη λειτουργία του εντέρου -- στόχος: μαλακές κενώσεις χωρίς προσπάθεια}
% \DrugEntry{tb Xanax}{0,25mg}{1 x 2-3}{}{μέχρι 3 την ημέρα επί εκνευρισμού}
% \DrugEntry{tb Escitalopram}{10mg}{1 x 1}{}{}
% \DrugEntry{spray Vioplex-T ή Pulvo47}{}{1 x 2}{}{στο τραύμα, και να καλύπτεται μετά}
% ---------------------------------------------------------------------
\end{MedTable}

% Notes under the table (uncomment the ones you need)
\textbf{Σημείωση:} Μην λαμβάνετε αντιφλεγμονώδη φάρμακα (π.χ.\ Voltaren, Brufen, Ponstan, Naprosyn) χωρίς ιατρική οδηγία.
% \textbf{Σημείωση:} Μην διακόψετε και μην μειώσετε μόνος/μόνη σας τα Τιμή 10. Οι δόσεις θα μειωθούν σταδιακά κατόπιν ιατρικής οδηγίας.
% \textbf{Σημείωση:} Ενημερώστε κάθε ιατρό και φαρμακοποιό ότι λαμβάνετε Τιμή 11, πριν από τη χορήγηση οποιουδήποτε νέου φαρμάκου.
% \textbf{Σημείωση:} Η υπόλοιπη φαρμακευτική σας αγωγή (π.χ.\ αντιυπερτασική) συνεχίζεται \textbf{αμετάβλητη}.

\end{document}
